% Part: lambda-calculus
% Chapter: syntax
% Section: alpha

\documentclass[../../../include/open-logic-section]{subfiles}

\begin{document}

\olfileid[pt]{lam}{syn}{alp}
\olsection{Conversão~$\alpha$}

Qual é a relação entre $\lambd[x][x]$ e $\lambd[y][y]$? Ambos representam a
função identidade. Naturalmente, são termos sintaticamente distintos. Eles
diferem apenas no nome da variável ligada, e um resulta de «renomear» a
variável ligada do outro. Isso é chamado de \emph{conversão~$\alpha$}.

\begin{defn}[Mudança de variável ligada, $\aconvone$]
  Se um termo $M$ contém uma ocorrência de $\lambd[x][N]$, $y \notin
  \FV{N}$, e $\Subst{N}{y}{x}$ está definida, então substituir essa ocorrência
  por
  \begin{equation*}
    \lambd[y][\Subst{N}{y}{x}]
  \end{equation*}
  dando como resultado $M'$ é chamado de \emph{mudança de variável ligada},
  escrita como $M \redone[\alpha] M'$.
\end{defn}

\begin{defn}[Compatibilidade de uma relação]
  Uma relação $R$ sobre termos é dita \emph{compatível} se satisfaz as
  seguintes condições:
  \begin{enumerate}
  \item Se $R N N'$, então $R \lambd[x][N] \lambd[x][N']$
  \item Se $R P P'$, então $R (PQ) (P'Q)$
  \item Se $R Q Q'$, então $R (PQ) (PQ')$
  \end{enumerate}
\end{defn}

Reformulemos, então, a definição:
\begin{defn}[Mudança de variável ligada, $\aconvone$]
  A \emph{mudança de variável ligada} ($\redone[\alpha]$) é a menor relação
  compatível sobre termos que satisfaz a seguinte condição:
  \begin{align*}
    &\lambd[x][N] \redone[\alpha] \lambd[y][\Subst{N}{y}{x}] && \text{se
      $x \neq y$, $y \notin \FV{N}$} \\
    & &&\text{e $\Subst{N}{y}{x}$ está definida}
  \end{align*}
\end{defn}

Aqui, «menor» significa que a relação contém apenas os pares exigidos pela
compatibilidade e pela condição adicional, e nada mais. Assim, essa relação
também pode ser definida da seguinte maneira:

\begin{defn}[Mudança de variável ligada, $\aconvone$] \ollabel{defn:aconvone}
  A \emph{mudança de variável ligada} ($\aconvone$) é definida indutivamente
  da seguinte maneira:
  \begin{enumerate}
  \item Se $N \aconvone N'$, então $\lambd[x][N] \aconvone
    \lambd[x][N']$ \ollabel{defn:aconvone1}
  \item Se $P \aconvone P'$, então $(PQ) \aconvone (P'Q)$ \ollabel{defn:aconvone2}
  \item Se $Q \aconvone Q'$, então $(PQ) \aconvone (PQ')$ \ollabel{defn:aconvone3}
  \item Se $x \neq y$, $y \notin \FV{N}$ e $\Subst{N}{y}{x}$ está definida,
    então $\lambd[x][N] \redone[\alpha] \lambd[y][\Subst{N}{y}{x}]$.
    \ollabel{defn:aconvone4}
  \end{enumerate}
\end{defn}

As definições são equivalentes, mas deixamos a demonstração como exercício.
Daqui em diante, usaremos a definição indutiva.

\begin{defn}[Conversão~$\alpha$, $\aconv$]
  A \emph{conversão~$\alpha$} ($\aconv$) é a menor relação reflexiva e
  transitiva sobre termos que contém~$\aconvone$.
\end{defn}

Como acima, «menor» significa que a relação contém apenas os pares exigidos
pela transitividade e por~$\aconvone$, o que conduz à seguinte definição
equivalente:

\begin{defn}[Conversão~$\alpha$, $\aconv$] \ollabel{defn:aconv}
  A \emph{conversão~$\alpha$} ($\aconv$) é definida indutivamente da seguinte
  maneira:
  \begin{enumerate}
  \item Se $P \aconv Q$ e $Q \aconv R$, então $P \aconv R$.
    \ollabel{defn:aconv1}
  \item Se $P \aconvone Q$, então $P \aconv Q$. \ollabel{defn:aconv2}
  \item $P \aconv P$. \ollabel{defn:aconv3}
  \end{enumerate}
\end{defn}

\begin{ex}
  $\lambd[x][f x]$ converte-se por~$\alpha$ em $\lambd[y][f y]$, e
  reciprocamente. Em termos informais, ambas são funções que recebem um
  argumento e devolvem $f$ aplicado a esse argumento, referindo-se à variável
  de ambiente~$f$.

  $\lambd[x][f x]$ não se converte por~$\alpha$ em $\lambd[x][g x]$. Em
  termos informais, elas se referem, respectivamente, às variáveis de ambiente
  $f$ e~$g$, o que as torna funções diferentes: elas se comportam de maneira
  distinta em ambientes nos quais $f$ e $g$ são diferentes.
\end{ex}

\begin{prob}
  Os pares de termos a seguir são conversíveis por~$\alpha$?
  \begin{enumerate}
  \item $\lambd[x][\lambd[y][x]]$ e $\lambd[y][\lambd[x][y]]$
  \item $\lambd[x][\lambd[y][x]]$ e $\lambd[c][\lambd[b][a]]$
  \item $\lambd[x][\lambd[y][x]]$ e $\lambd[c][\lambd[b][a]]$
  \end{enumerate}
\end{prob}

\begin{lem}\ollabel{lem:fv-one}
  Se $P \aconvone Q$, então $\FV{P} = \FV{Q}$.
\end{lem}

\begin{proof}
  Por indução sobre a !!{derivation} de $P \aconvone Q$.
  \begin{enumerate}
  \item Se a última regra é \olref{defn:aconvone4}, então $P$ tem a forma
    $\lambd[x][N]$ e $Q$ tem a forma $\lambd[y][\Subst{N}{y}{x}]$, com
    $x \neq y$, $y \notin \FV{N}$ e $\Subst{N}{y}{x}$ definida. Distinguimos
    casos conforme $x \in \FV{N}$ ou não:
    \begin{enumerate}
    \item Se $x \in FV(N)$, então:
      \begin{align*}
        \FV{\lambd[y][\Subst{N}{y}{x}]} &=
        \FV{\Subst{N}{y}{x}} \setminus \{y\} \\
        & = ((\FV{N} \setminus \{x\}) \cup \{y\}) \setminus \{y\}
         && \text{por \olref[sub]{thm:infv}} \\
        & = \FV{N} \setminus \{x\} \\
        & = \FV{\lambd[x][N]}
      \end{align*}
    \item Se $x \notin FV(N)$, então:
      \begin{align*}
        \FV{\lambd[y][\Subst{N}{y}{x}]}
        & = FV{\Subst{N}{y}{x}} \setminus \{y\} \\
        & = \FV{N} \setminus \{x\}
         && \text{por \olref[sub]{thm:notinfv}} \\
        & = \FV{\lambd[x][N]}.
      \end{align*}
    \end{enumerate}
  \item Os outros três casos ficam como exercícios.
  \end{enumerate}
\end{proof}

\begin{prob}
  Complete a demonstração de \olref[lam][syn][alp]{lem:fv-one}.
\end{prob}

\begin{lem}\ollabel{lem:inv}
  Se $P \aconvone Q$, então $Q \aconvone P$.
\end{lem}

\begin{proof}
  Indução sobre a !!{derivation} de $P \aconvone Q$.
  \begin{enumerate}
  \item Se a última regra é \olref{defn:aconvone4}, então $P$ tem a forma
    $\lambd[x][N]$ e $Q$ tem a forma $\lambd[y][\Subst{N}{y}{x}]$, onde
    $x \neq y$, $y \notin \FV{N}$ e $\Subst{N}{y}{x}$ está definida.
    Primeiro, temos $y \notin \FV{\Subst{N}{y}{x}}$ por
    \olref[sub]{thm:clr}. Por \olref[sub]{thm:inv}, temos que
    $\Subst{\Subst{N}{y}{x}}{x}{y}$ não apenas está definida, mas também é
    igual a~$N$. Então, por \olref{defn:aconvone4}, temos
    $\lambd[y][\Subst{N}{y}{x}] \aconvone
    \lambd[x][\Subst{\Subst{N}{y}{x}}{x}{y}] = \lambd[x][N]$.
  \end{enumerate}
\end{proof}

\begin{prob}
  Complete a demonstração de \olref[lam][syn][alp]{lem:inv}.
\end{prob}

\begin{thm}
  A conversão~$\alpha$ é uma relação de equivalência sobre termos; O.s.\@,
  ela é reflexiva, simétrica e transitiva.
\end{thm}

\begin{proof}
  \begin{enumerate}
  \item Para cada termo $M$, $M$ pode ser mudado para $M$ por
    \emph{zero} mudanças de variáveis ligadas.
  \item Se $P$ se converte por~$\alpha$ em $Q$ mediante uma sequência de
    mudanças de variáveis ligadas, então, a partir de $Q$, podemos simplesmente
    inverter essas mudanças (por \olref{lem:inv}) na ordem oposta para
    obter~$P$.
  \item Se $P$ se converte por~$\alpha$ em $Q$ mediante uma sequência de
    mudanças de variáveis ligadas, e $Q$ em $R$ mediante outra sequência,
    então podemos mudar $P$ para $R$ aplicando primeiro a primeira sequência e
    depois a segunda.
  \end{enumerate}
\end{proof}

Daqui em diante, dizemos que $M$ e $N$ são \emph{$\alpha$-equivalentes},
$M \aeq N$, se e somente se $M$ se converte por~$\alpha$ em $N$ (o que, como
acabamos de mostrar, ocorre se e somente se $N$ se converte por~$\alpha$
em~$M$).

\begin{thm}\ollabel{thm:fv}
  Se $M \aeq N$, então $\FV{M} = \FV{N}$.
\end{thm}

\begin{proof}
  Imediato de \olref{lem:fv-one}.
\end{proof}

\begin{lem}\ollabel{lem:sub:R}
  Se $R \aeq R'$ e $\Subst{M}{R}{y}$ está definida, então
  $\Subst{M}{R'}{y}$ está definida e é $\alpha$-equivalente a
  $\Subst{M}{R}{y}$.
\end{lem}

\begin{proof}
  Exercício.
\end{proof}

\begin{prob}
  Demonstre \olref[lam][syn][alp]{lem:sub:R}.
\end{prob}

Lembre-se de que, em \olref[sub]{sec}, a substituição não está definida em
alguns casos; contudo, usando a conversão~$\alpha$ sobre termos, podemos fazer
com que a substituição esteja sempre definida por meio da renomeação de
variáveis ligadas. O resultado preserva a $\alpha$-equivalência, como mostra o
teorema a seguir:

\begin{thm}\ollabel{thm:sub}
  Para quaisquer $M$, $R$ e~$y$, existe $M'$ tal que $M \aeq M'$ e
  $\Subst{M'}{R}{y}$ está definida. Além disso, se houver outro par
  $M'' \aeq M$ e $R''$ para o qual $\Subst{M''}{R''}{y}$ esteja definida e
  $R'' \aeq R$, então $\Subst{M'}{R}{y} \aeq \Subst{M''}{R''}{y}$.
\end{thm}

\begin{proof}
  Por indução sobre a formação de $M$:
  \begin{enumerate}
  \item $M$ é uma variável~$z$: exercício.
  \item Suponha que $M$ tenha a forma $\lambd[x][N]$. Escolha uma variável
    $z$ distinta de $x$ e $y$ e tal que $z \notin \FV{N}$ e
    $z \notin \FV{R}$. Pela hipótese de indução, existe $N'$ tal que
    $N' \aeq N$ e $\Subst{N'}{z}{x}$ está definida. Então
    $\lambd[x][N] \aeq \lambd[x][N']$ também, por
    \olref{defn:aconvone}\olref{defn:aconvone1}. Agora,
    $\lambd[x][N'] \aeq \lambd[z][\Subst{N'}{z}{x}]$ por
    \olref{defn:aconvone}\olref{defn:aconvone4}. Podemos fazer isso porque
    $z \ne x$, $z \notin FV(N')$ e $\Subst{N'}{z}{x}$ está definida.
    Finalmente, $\Subst{\lambd[z][\Subst{N'}{z}{x}]}{R}{y}$ está definida,
    porque $z \neq y$ e $z \notin FV(R)$.

    Além disso, se houver outros $N''$ e $R''$ que satisfaçam as mesmas
    condições,
    \begin{multline*}
      \Subst{(\lambd[z][\Subst{N''}{z}{x}])}{R''}{y} =\\
    \begin{aligned}
      &= \lambd[z][\Subst{\Subst{N''}{z}{x}}{R''}{y}] \\
      &= \lambd[z][\Subst{\Subst{N''}{z}{x}}{R}{y}] && \text{por
        \olref{lem:sub:R}}\\
      &=\lambd[z][\Subst{\Subst{N'}{z}{x}}{R}{y}]
      && \text{pela hipótese de
        indução}\\
      &=\Subst{(\lambd[z][\Subst{N'}{z}{x}])}{R}{y}
    \end{aligned}
    \end{multline*}
    \item $M$ é da forma $(PQ)$: exercício.
  \end{enumerate}
\end{proof}

\begin{prob}
  Complete a demonstração de \olref[lam][syn][alp]{thm:sub}.
\end{prob}

\begin{cor}\ollabel{cor:sub}
  Para quaisquer $M$, $R$ e~$y$, existe um par $M'$ e $R'$ tal que
  $M' \aeq M$, $R \aeq R'$ e $\Subst{M'}{R'}{y}$ está definida. Além disso,
  se houver outro par $M'' \aeq M$ e $R''$ com $\Subst{M'}{R'}{y}$ definida,
  então $\Subst{M'}{R'}{y} \aeq \Subst{M''}{R''}{y}$.
\end{cor}

\begin{proof}
  Imediato de \olref{thm:sub}.
\end{proof}

\end{document}
